\section{Freedom, love, and the gift of blessedness}
\label{sec:angel-grace}

\subsection{Freedom without indecision (Question 59)}

A creature may know what is good without being identical with goodness. It can also love a good without producing it by that love. Aquinas therefore places will in the angel as an intellectual appetite: an inclination towards the good as understood. ``Appetite'' does not necessarily mean hunger, lack, or restlessness. It includes delight in a good already possessed. Blessedness can perfect will without making willing cease.\footnote{\ST{I q.~59 a.~1}.}

Intellect and will differ by their relations to their objects. Knowing receives the object's intelligible form within the knower; willing tends towards the good of the object. This explains why an angel's will is neither its intellect nor its essence. An angel can understand evil, but can desire it only under some appearance or aspect of good. God alone contains universal goodness in his own essence.\footnote{\ST{I q.~59 a.~2}.}

Freedom consequently does not require the delay of wondering what to do. Article 3 distinguishes the discursive deliberation characteristic of human choice from the free judgment belonging to intellectual nature. An angel can choose without an interval of inquiry. Certainty about a thing's nature also need not make every possible action towards it necessary. The good can be pursued through diverse acts, and created freedom has a genuine field of election.\footnote{\ST{I q.~59 a.~3}, corpus and ad~1--2.}

Article 4 denies angels the bodily sensitive appetites called concupiscible and irascible. Yet it does not empty Scripture's descriptions of love and joy. Those can name acts of will without naming bodily passions. Courage can likewise designate steadfast execution of God's will, though there is no bodily fear to moderate. The vocabulary remains meaningful, with a different mode of realization.\footnote{\ST{I q.~59 a.~4}.}

\angeldossiersettled{I q.~59, aa.~1--4}{Church teaching: angels possess will (\emph{CCC} 330). Common teaching: their personal freedom. Thomist position: the distinction of faculties and the analysis of affection without sensitive appetite.}{Augustine, \emph{Trinity} X; Aristotle, \emph{De anima} III; Dionysius, \emph{Divine Names} IV, as discussed by Aquinas.}

\subsection{Loving oneself within a greater good (Question 60)}

Natural love is not a second will alongside free will. It is the creature's foundational inclination to its own good and perfection. Elective love concerns goods chosen in relation to an end. Aquinas compares the relation with intellect's first principles and the conclusions resting on them: choice does not create its own first orientation towards goodness.\footnote{\ST{I q.~60 aa.~1--2}.}

Thus self-love can be natural or elective. An angel naturally seeks its good; it also chooses particular goods for itself. Loving another means willing that other's good, not merely wishing to possess or use the other. Article 4 grounds natural mutual love in the unity angels share as intellectual creatures. ``As oneself'' signifies a likeness in benevolence, not equality of intensity or identity of will. Even fallen angels retain natural affinity with good angels while hating their righteousness. Nature's goodness and chosen moral opposition are different considerations.\footnote{\ST{I q.~60 aa.~3--4}, especially a.~4 ad~2--3.}

Article 5 reaches the difficult question: does natural love place God before self? Aquinas argues that an individual belongs to a good greater than its private good. A hand instinctively protects the whole body. The analogy does not make God a whole assembled from creaturely parts; it illustrates an inclination towards the good on which one's own good depends. Every creature receives all its being and goodness from God. Natural love therefore tends first towards God as universal good. If charity had to destroy that orientation, nature would already be perverse before sin.\footnote{\ST{I q.~60 a.~5}, corpus and ad~1--3.}

Natural love is nevertheless not charity. God as source of created goodness and God as the supernatural good enjoyed in the vision of his essence are distinguishable aspects of the same God. The second exceeds every created nature. The question therefore prepares the account of grace rather than making grace unnecessary.\footnote{\ST{I q.~60 a.~5 ad~4--5}.}

\angeldossier{I q.~60, aa.~1--5}{Thomist position: the natural and elective analysis, including natural love of God above self. Common teaching: charity is a supernatural gift.}{Augustine, \emph{Trinity} X and \emph{City of God} XII.9; Aristotle, \emph{Ethics} IX; Ecclesiasticus~13:19.}{In a.~5 Aquinas disputes a natural self-love held to be absolutely stronger than love of God, while distinguishing natural love from charity.}

\subsection{A beginning received from God (Question 61)}

An angel does not emerge from pre-existing spiritual material. Since its essence is not its existence, its whole being must be given by God. No creature shares the Creator's unoriginated eternity. Natural incorruptibility tells us that an angel does not decay; it cannot tell us that the angel never began.\footnote{\ST{I q.~61 aa.~1--2}; Lateran IV, \emph{Firmiter}, DS~800; \emph{CCC} 327--328.}

The relation of that beginning to the bodily world receives a more qualified answer. Aquinas judges simultaneous creation more probable because angels and bodies belong to one ordered universe. He expressly permits the contrary patristic position. Gregory Nazianzen describes the intellectual creation first, then the material world; John Damascene follows him. Aquinas's respect for that alternative prevents a theological preference from becoming a condemnation. Both positions confess the single Creator of spiritual and bodily things.\footnote{\ST{I q.~61 a.~3}; Gregory Nazianzen, Oration~38.9--10; Damascene, \emph{Orthodox Faith} II.3.}

Article 4 locates angelic creation fittingly in the empyrean, the highest bodily heaven of the received cosmology. Its central argument concerns their relation to the material world they serve. The place does not cause or sustain their existence: Aquinas expressly says God could have created angels before every body. The place expresses their ordered relation to bodies, while their being comes wholly from God.\footnote{\ST{I q.~61 a.~4}, especially ad~1.}

\angeldossier{I q.~61, aa.~1--4}{Defined: creation by God and a beginning of the created order (Lateran IV, DS~800). Disputed: relative sequence and the empyrean account in aa.~3--4.}{Psalm~148:2, 5; Proverbs~8:22; Genesis~1:1; Gregory Nazianzen, Oration~38.9--10; Damascene II.3.}{Gregory and Damascene place angelic creation before the material world. Aquinas favors simultaneity but explicitly declines to condemn the alternative.}

\subsection{Grace before merit, glory after charity (Question 62)}

A child may receive everything needed for a human life without receiving everything that child will ever become. The comparison is limited, because an angel does not mature by bodily growth; it nonetheless helps distinguish an initial nature from a final perfection. Aquinas says an angel begins with its natural intellectual perfection, but not with the supernatural vision of God. Article 1 therefore distinguishes natural happiness from consummated beatitude. The fallen angels never possessed the latter.\footnote{\ST{I q.~62 a.~1}; \ST{I q.~64 a.~1 ad~4}.}

Grace is necessary even before sin. It does more than repair a damaged faculty: it raises a creature to a life above its natural capacity. Article 2 distinguishes God as author of nature from God as the object of supernatural beatitude. The absence of bodily obstacles cannot bridge that distance. Article 3 favors creation in sanctifying grace, with Augustine's account of God creating the good angels' nature and holy love together. Aquinas judges this the more probable opinion. Grace is the beginning of a free good act; it does not force the creature to use it.\footnote{\ST{I q.~62 aa.~2--3}, especially a.~3 ad~2; Augustine, \emph{City of God} XII.9.}

Augustine's argument reaches beyond the origin of an individual good choice. If the angels had made their own will good without God's gift, they would have made themselves better than their Creator made them. Their holy love must itself be received. Through that love they cleave both to God and to one another: the city of God begins as a communion of blessed spirits, and redeemed humanity is gathered into its fellowship. Grace establishes a common life, not merely a private elevation.\footnote{Augustine, \emph{City of God} XII.9.}

Merit follows that gift. The angel does not manufacture the vision of God by natural achievement. Through charity it freely acts towards a reward that God bestows. In articles 4--5 Aquinas places one meritorious act before immediate beatification. ``Before'' signifies the ordering and succession of acts, not a measured probation lasting so many human hours. Article 6 then proposes that the grades of grace and glory correspond to natural excellence. His reason is God's wise and free ordering of both gifts. He explicitly distinguishes this angelic proposal from the varied relation of natural ability and grace in human beings.\footnote{\ST{I q.~62 aa.~4--6}, especially a.~5 ad~2 and a.~6 ad~1, 3.}

The final articles describe fulfillment. Nature's knowledge and love remain within glory (a.~7). The blessed angel cannot sin because it sees the divine goodness itself; freedom achieves its end without retaining the defect of choosing against that end (a.~8). Essential beatitude does not increase beyond the measure given by God, yet joy in the salvation of others can increase (a.~9). Service therefore expresses already attained perfection. It is neither a new competition for rank nor a suspension of happiness.\footnote{\ST{I q.~62 aa.~7--9}, especially a.~9 ad~2--3.}

\angeldossier{I q.~62, aa.~1--9}{Common teaching: grace is needed for supernatural beatitude and blessed angels are confirmed in good. Thomist position: one meritorious act, immediate reward, and correspondence of natural and gracious grades. Disputed: creation in grace is expressly judged more probable in a.~3.}{Romans~6:23; Augustine, \emph{City of God} XII.9; Peter Lombard, \emph{Sentences} II d.~3, invoked in a.~6; Matthew~18:10; Luke~15:10.}{Aquinas records creation in nature alone before grace, and rejects merit continuing to earn essential beatitude after that beatitude is possessed.}
